\section{Software and Compute}
\label{app:compute}

\paragraph{Models.}
We use the instruction-tuned checkpoints from Hugging Face: \texttt{Qwen2.5-0.5B-Instruct}, \texttt{Qwen2.5-7B-Instruct} and \texttt{Qwen2.5-14B-Instruct} \citep{qwen2025qwen25}, \texttt{OLMo-2-1124-7B-Instruct} \citep{olmo2024olmo2}, and \texttt{granite-3.1-8b-instruct} \citep{granite2024}. We do not fine-tune any model.

\paragraph{Running the models.}
All models run in 32-bit floating point with Hugging Face Transformers and PyTorch, on CPU nodes of a university cluster. Each model answers all 1,712 abstention prompts (856 questions, closed book and with the page) in one 12-thread job. We score prompts in batches with left padding. We rebuild the position indices from the attention mask, so a batched prompt gets exactly the same probabilities as the same prompt alone. Every answer is stored under a hash of its exact prompt, so an edited page never reuses an old answer.

\paragraph{Recovery.}
In recovery, all candidates for one fact share the same prompt prefix. We encode this prefix once and reuse its key--value cache for every candidate. Before each run, we check a sample of these cached scores against full forward passes.

\paragraph{Randomness and statistics.}
All sampling uses seed 0, and the bootstrap draws are sorted by question identifier before resampling, so every interval in the paper is the same on every run. The statistics use only the Python standard library.
